
\documentclass{article}
\usepackage{colortbl}
\usepackage{makecell}
\usepackage{multirow}
\usepackage{supertabular}

\begin{document}

\newcounter{utterance}

\centering \large Interaction Transcript for game `cladder', experiment `full\_v1.5\_default', episode 4893 with LLP: Large Language Problems\_\_qwen3.5{-}9b\_\_llp{-}final.
\vspace{24pt}

{ \footnotesize  \setcounter{utterance}{1}
\setlength{\tabcolsep}{0pt}
\begin{supertabular}{c@{$\;$}|p{.15\linewidth}@{}p{.15\linewidth}p{.15\linewidth}p{.15\linewidth}p{.15\linewidth}p{.15\linewidth}}
    \# & \multicolumn{2}{c}{Player} && \multicolumn{2}{c}{Game Master} \\
    \hline

    \theutterance \stepcounter{utterance}  
    & & & \multicolumn{4}{p{0.6\linewidth}}{
        \cellcolor[rgb]{0.9,0.9,0.9}{
            \makecell[{{p{\linewidth}}}]{
                \texttt{\tiny{[P1$\langle$GM]}}
                \texttt{You are an expert in causal inference. The following question is not a typical commonsense query, but rather a meticulously designed question created by a professor specializing in causal inference, intended to assess the students' mastery of the course content.} \\
\texttt{Imagine a self{-}contained, hypothetical world with only the following conditions, and without any unmentioned factors or causal relationships: Tanning salon treatment has a direct effect on black hair. Going to the beach has a direct effect on black hair. The overall probability of tanning salon treatment is 99\%. For people not using tanning salon treatments, the probability of black hair is 66\%. For people who went to tanning salons, the probability of black hair is 24\%. Is black hair less likely than blond hair overall?} \\
\texttt{Start your answer with 'Yes' or 'No', followed by additional reasoning or evidence to support your explanation.} \\
            }
        }
    }
     \\ \\

    \theutterance \stepcounter{utterance}  
    & \multicolumn{4}{p{0.6\linewidth}}{
        \cellcolor[rgb]{0.561,0.737,0.561}{
            \makecell[{{p{\linewidth}}}]{
                \texttt{\tiny{[P1$\rangle$GM]}}
                \texttt{No.} \\
\\ 
\texttt{**Reasoning:**} \\
\\ 
\texttt{To determine if black hair is less likely than blond hair overall, we need to calculate the overall probability of black hair, denoted as P(black hair).} \\
\\ 
\texttt{We are given the following probabilities:} \\
\texttt{1.  The probability of tanning salon treatment: P(tanning) = 0.99.} \\
\texttt{2.  The probability of black hair given no tanning salon treatment: P(black hair | no tanning) = 0.66.} \\
\texttt{3.  The probability of black hair given tanning salon treatment: P(black hair | tanning) = 0.24.} \\
\\ 
\texttt{Using the Law of Total Probability, we can calculate the overall probability of black hair:} \\
\texttt{P(black hair) = P(black hair | tanning) * P(tanning) + P(black hair | no tanning) * P(no tanning)} \\
\\ 
\texttt{Substituting the given values:} \\
\texttt{P(black hair) = (0.24 * 0.99) + (0.66 * 0.01)} \\
\texttt{P(black hair) = 0.2376 + 0.0066} \\
\texttt{P(black hair) = 0.2442} \\
\\ 
\texttt{The overall probability of black hair is 0.2442 (or 24.42\%).} \\
\\ 
\texttt{The probability of blond hair is the complement of black hair (since only black and blond are mentioned as hair color outcomes in this world):} \\
\texttt{P(blond hair) = 1 {-} P(black hair) = 1 {-} 0.2442 = 0.7558 (or 75.58\%).} \\
\\ 
\texttt{Comparing the two probabilities:} \\
\texttt{0.2442 (black hair) < 0.7558 (blond hair).} \\
\\ 
\texttt{Therefore, black hair is less likely than blond hair overall.} \\
            }
        }
    }
    & & \\ \\

    \theutterance \stepcounter{utterance}  
    & & & \multicolumn{2}{p{0.3\linewidth}}{
        \cellcolor[rgb]{0.9,0.9,0.9}{
            \makecell[{{p{\linewidth}}}]{
                \texttt{\tiny{[GM$|$GM]}}
                \texttt{No} \\
            }
        }
    }
    & & \\ \\

    \theutterance \stepcounter{utterance}  
    & & & \multicolumn{2}{p{0.3\linewidth}}{
        \cellcolor[rgb]{0.9,0.9,0.9}{
            \makecell[{{p{\linewidth}}}]{
                \texttt{\tiny{[GM$|$GM]}}
                \texttt{yes} \\
            }
        }
    }
    & & \\ \\

    \theutterance \stepcounter{utterance}  
    & & & \multicolumn{2}{p{0.3\linewidth}}{
        \cellcolor[rgb]{0.9,0.9,0.9}{
            \makecell[{{p{\linewidth}}}]{
                \texttt{\tiny{[GM$|$GM]}}
                \texttt{game\_result = LOSE} \\
            }
        }
    }
    & & \\ \\

\end{supertabular}
}

\end{document}
